\section{Conclusions}
\label{sec:conclusions}

This paper extends the shape-based framework of Papers~I--III from the scalar
mean surface gravity to the spatially resolved surface fields---effective
gravity, gravitational slope, and geopotential---computed facet by facet from a
polyhedral shape model and an assumed uniform density. The main conclusions are
as follows.

\begin{enumerate}

\item \emph{The implementation reproduces an independent calculation facet by
facet.} Compared against the gravitational products released with the Hayabusa2
shape analysis of Ryugu, on the identical $49\,152$-facet model at the same
density and rotation period, the area-weighted mean slope agrees to better than
$0.01^{\circ}$, the RMS per-facet difference is $0.48^{\circ}$, $99.96\%$ of
facets agree to within $5^{\circ}$ and $99.95\%$ to within $2^{\circ}$, and the
Pearson
correlations are $0.9989$ in slope and $0.9898$ in geopotential. The agreement
holds band by band in latitude to better than $0.7^{\circ}$.

\item \emph{From a shape model, an assumed density and a specified spin state
the framework recovers the published surface gravity structure of Bennu.} The global mean slope is
$15.45^{\circ}$; the latitudinal profile rises from $10.45^{\circ}$ near the
equator to $19.89^{\circ}$ at mid-latitude; the geopotential minimum lies at the
equator, as inferred from mission data. Sweeping the assumed bulk density
reproduces the published slope--density relation quantitatively, returning
$25.55^{\circ}$ and $16.06^{\circ}$ at $0.85$ and $1.15\ \mathrm{g\,cm^{-3}}$
against roughly $24^{\circ}$ and $15^{\circ}$ reported from OSIRIS-REx analyses.

\item \emph{The area-weighted mean slope did not converge within the tested
mesh range on shape models that carry substantial unresolved relief.} Across
meshes from about $2\,000$ to $42\,000$ facets it varies by $1.8\%$ to $25.6\%$
depending on the body, and on three of five bodies it was still moving by
$2.6$--$8.2\%$ over the final refinement step, with no plateau. Laplacian
smoothing at fixed facet count lowers the slope by an amount comparable to what
refinement adds, identifying surface relief as the cause. A mean slope is
therefore a statement about a particular mesh and is not reproducible without
its facet count.

\item \emph{The normalized geopotential dispersion, computed from the same
field on the same meshes, shows clear practical convergence over the same
range.} It varies by $0.9\%$ to $5.9\%$ across the
same range and changes by at most $0.60\%$ over the final refinement step---by
nothing at all on two bodies---and is the more stable of the two diagnostics on
every body tested. The distinction follows from what each quantity depends on:
the geopotential is an integral over the mass distribution and inherits the
convergence of the polyhedron's volume, whereas the slope depends on the
orientation of an individual facet, which a rough surface continues to change
at every scale at which it is resolved.

\item \emph{In the controlled Itokawa comparison the ground-based model
preserves the global gravity descriptors but not the local slope statistics.}
Comparing the spacecraft-derived and radar-derived models at approximately
matched resolution, the area-weighted mean slope agrees to $0.8\%$ and the
geopotential dispersion to $7.9\%$, while the radar model contains no facet
steeper than $41^{\circ}$ against $149^{\circ}$ for the spacecraft model, and
places $1.7\%$ rather than $4.8\%$ of the surface above $30^{\circ}$. On the
evidence of this single controlled case, population-level characterisation from
ground-based shapes appears defensible, whereas identification of steep terrain,
overhangs, or the fraction of surface above the angle of repose does not.

\item \emph{A rotating test case is necessary to verify a small-body gravity
code.} A non-rotating sphere returns zero slope for either sign of the
centrifugal term and so cannot detect an error in the rotating frame. The
uniformly rotating sphere, whose slope field is given in closed form by
Eq.~\eqref{eq:rotsphere} and derived in Appendix~\ref{app:rotsphere},
discriminates sharply---the two signs differ by more than seven degrees at
mid-latitude for $\omegarot^{2}R/g = 0.5$. We recommend it as a routine check
and provide the implementation.

\end{enumerate}

Two qualifications frame these results. For bodies without a measured mass the
assumed bulk density remains the dominant uncertainty, shifting the mean slope
by some $10^{\circ}$ on Bennu across a plausible range---three times the effect
of mesh resolution---so that constraining the mass buys more than refining the
shape. And all fields assume a homogeneous interior, an assumption that leaves
the resolution results untouched but that is known to fail for bilobed bodies
such as Itokawa, whose head and body differ appreciably in density.

That last limitation points to the natural continuation. A surface shaped by
long-term downslope transport should approach an equipotential of the true
gravity field, so that the spread of geopotential over the surface becomes an
estimator of the interior density distribution rather than merely a description
of it. Such an inversion requires a diagnostic whose response to the assumed
density is not swamped by its response to the mesh---the property established
here for the geopotential dispersion and shown to be absent for the mean slope.
That development is reserved for future work and lies beyond the scope of the
present study.
% ============================================================
% DATA AVAILABILITY STATEMENT  --  SDGA Paper IV
%
% Paste this block at the END of Section8.tex, after the final
% paragraph of the Conclusions and before \end{document} is reached.
%
% The reserved DOI is already filled in below (10.5281/zenodo.21793575).
% No further editing is needed.
% ============================================================

\section*{Data availability}

The analysis code, the verification scripts, the five figures, and a
machine-readable file containing every numerical result quoted in this paper
are openly available at
\href{https://doi.org/10.5281/zenodo.21793575}{doi:10.5281/zenodo.21793575},
released under CC BY 4.0. The deposit includes the two verification programs
that implement Eqs.~\eqref{eq:ellipsoid} and~\eqref{eq:rotsphere}, so that both
checks can be reproduced independently of the shape models. Development of the
code is tracked at
\url{https://github.com/praponkan/shape-based-surface-gravity}.

Shape models are not redistributed. Those used here are available from the
OSIRIS-REx mission archive (Bennu), the JAXA DARTS archive accompanying
\citet{kanamaru2019} (Ryugu shape model and reference gravitational products),
and the PDS Small Bodies Node (Eros, Itokawa, Gaspra, Ida, Kleopatra); the
specific models, assumed densities, and rotation periods are listed in
Table~\ref{tab:shapes}.


% ============================================================
% OPTIONAL -- if the journal requires a separate acknowledgement
% of the software used, this can follow the statement above.
% ============================================================

% \section*{Software}
%
% Self-gravitation was computed with \texttt{polyhedral-gravity}
% version~3.3.1 \citep{gravitylib}. Analysis used Python~3.13 with
% \texttt{numpy} and \texttt{scipy}; figures were produced with
% \texttt{matplotlib}.
